\subsection{NIL Prediction}

\begin{definition}[NIL Prediction]
    Given an entity mention $m$ and a reference \ac{kr}, \emph{NIL prediction} is the subtask of \acl{nel} that identifies whether $m$ referes to an entity that does not exist in the \ac{kr}~\cite{entity_linking_with_kb_2015}. 
\end{definition}

According to some work \emph{NIL} stands for ``not in lexicon''~\cite{ILIEVSKI2020100565}. Another interpretation is that it simply refers to the English word ``nil'' that means ``nothing''~\cite{cambridge_nil}.

This task has also been called with different names: ``unlinkable mention prediction''~\cite{entity_linking_with_kb_2015}, ``out-of-knowledge-Base mention discovery''~\cite{dong_reveal_the_unknown_2023}, ``out-of-knowledge-graph entity detection''~\cite{moller2024entity}, ``unknwon entity discovery''~\cite{kassner-etal-2022-edin}, and ``emerging entity discovery''~\cite{hoffart_discovering_emergin_entities_2014}.

\textcite{zhou-etal-2024-gendecider} distinguishes between NIL prediction and ``none of the candidates'' detection, arguing that \emph{NIL} means the referred entity is not in the \ac{kr}, while ``none of the candidates'' means the referred entity has not been retrieved without assuming it is not present in the \ac{kr}. However, in practise, NIL prediction is often performed by predicting whether the top-ranked entity, returned by the \ac{nel} ranking step, is correct or not~\cite{entity_linking_with_kb_2015}. It is impractical and prohibitive to compare a mention with every entity in a \ac{kr} for ensuring the mentions is NIL. Thus, in this work, we consider the preactical assumption that if a mention does not refer to the top-ranked candidate(s), returned by the \ac{nel} system, the mention is NIL. Under this assumption there is no difference between NIL prediction and none-of-the-candidates prediction.

NIL prediction has frequently been overlooked in \acrshort{nel}
systems, as many approaches exclude NIL mentions from their evaluation datasets and
eventually leave this task to future works. In fact, among the 38 approaches\todo{verify the published version}
compared by the survey \citet{Sevgili2022NeuralEL} only 8 addressed NIL prediction.

% dico che le TAC hanno dato importanza al NIL
The origins of the task can be traced to the TAC\footnote{\url{https://tac.nist.gov/}} (Text Analysis
Conference), which introduced NIL prediction in its \acrshort{kbp} track starting from
2009 \cite{mcnamee2009overview}. Starting from 2011, participants were also required to cluster
together those NIL mentions that referred to the same entity \cite{ji2011tac}. Organized by NIST\footnote{\url{https://www.nist.gov/}}, TAC aimed to advance research in \acrshort{nlp} and related areas by offering large-scale test collections and standardized evaluation procedures. Since 2008, it has been held annually, inspiring a variety of studies on NIL prediction.



\acrlong{nel} with NIL prediction is essentially a classification with a reject
option that consists in NIL or unlinkable. There are four common approaches to
perform NIL prediction \cite{sevgili2021neural}:
\begin{enumerate}[leftmargin=1.5cm,label=\textit{NIL\textsubscript{\arabic*}:},ref=\textit{NIL\textsubscript{\arabic*}}]
    \item \label{nilm-empty-set} Linking to NIL when the candidate generation
          gives an empty set of candidates.
    \item \label{nilm-threshold} Setting a threshold for the linking score,
          below which a mention is considered unlinkable.
    \item \label{nilm-add-class} Adding a NIL entity to the candidates set, in
          the ranking phase. Hence a new class is added to a classification
          problem.
    \item \label{nilm-binary-clf} Training an additional binary classifier that
          accepts as input mention-entity pairs possibly with additional
          features, such as best linking score or information from
          \acrshort{ner}, and that finally classifies whether a mention is
          linkable or not.
\end{enumerate}
Some publications \cite{Zhang2012TheNE, Rao2013EntityLF} state method
\ref{nilm-threshold} (threshold) and \ref{nilm-add-class} (additional class) are
equivalent even though the latter requires no hand-tuning to set the threshold.

Generally, these approaches come inside a pipeline-based system that considers NIL
prediction and \acrshort{nel} as separate tasks. Some works, in opposition,
proposed a single joint system to avoid error propagation and to exploit
inter-task dependencies \cite{fahrni2013hits,martins-etal-2019-joint}.


\subsection{NIL Clustering}
\subsection{coreference resolution}
\todo{what did LLMs change?}
\subsection{Human-in-the-Loop}
